\documentclass[11pt]{article}
\usepackage{amsmath,amssymb}
\title{Literal local action, gauge-Jacobian and Haar contact identities}\author{}\date{}
\begin{document}\maketitle
\paragraph{Target.}
By the short line contact before current decay calculation, the actual endpoint minor retains a genuine
short-line current contact, so a per-minor longest-time improvement is
false. Here the local original tensors are matched in the gauge sector.
This does not identify the surviving transverse contact of the short line contact before current decay calculation with a
longitudinal one.

Use all free links of the original relative cube and interior gauge
parameters zero on its complete boundary. At flat background write
\[
 S_0(\alpha z)=\alpha^2 S_2(z)+\alpha^3T_3(z)+\alpha^4T_4(z)+\cdots,
 \quad S_2=\tfrac12z^THz,\quad
 T_3=t_{abc}z_az_bz_c/6,\quad T_4=u_{abcd}z_az_bz_cz_d/24.
\]
The local logarithmic gauge velocity from the local gauge completion first nonlinear coefficient calculation is
\[
 \mathcal R_\xi(z)=D\xi+\alpha R_1(z)\xi+\alpha^2R_2(z)\xi+\cdots,
 \quad R_1=-\tfrac12\operatorname{ad}_z A,\quad
 R_2=\tfrac1{12}\operatorname{ad}_z^2D.                           \tag{1}
\]
Every tensor is the actual finite-lattice tensor, including links
normal to boundary faces; no periodic or infinite-volume rule is used.

\paragraph{1. Original cubic and quartic contacts.}
Gauge invariance of the ORIGINAL plaquette action gives, order by order,
\[
 HD=0,\qquad
 \partial_{D\xi}T_3+\langle Hz,R_1(z)\xi\rangle=0,                 \tag{2}
\]
\[
 \boxed{\partial_{D\xi}T_4+
       \partial_{R_1(z)\xi}T_3+\langle Hz,R_2(z)\xi\rangle=0.}    \tag{3}
\]
These include the nonlinear gauge response, not only a constant
gradient shift. To make the coefficients unambiguous put
$(R_1)_a=A_{abv}z_b\xi_v$ and
$(R_2)_a=B_{abcv}z_bz_c\xi_v/2$, with $B$ symmetric in $b,c$.
Then
\[
 D_{av}t_{abc}+H_{ab}A_{acv}+H_{ac}A_{abv}=0,                    \tag{4}
\]
\[
 \begin{split}
 D_{av}u_{abcd}
 &+t_{abc}A_{adv}+t_{abd}A_{acv}+t_{acd}A_{abv}\\
 &+H_{ab}B_{acdv}+H_{ac}B_{abdv}+H_{ad}B_{abcv}=0.               \end{split}\tag{5}
\]
Thus the quartic contraction, cubic response and second gauge-response
contact have fixed coefficients. Treating (5) as just an unnamed Ward
identity would hide precisely the partners needed in a contraction.

\paragraph{2. Collapse an actual longitudinal propagator leg.}
Let $G=(H+DD^*)^{-1}$ and $g=(D^*D)^{-1}$.
At the actual critical background, not only at flat background,
\[
 (H+DD^*)D=D(D^*D),\quad GD=Dg,\quad
 D^*e^{-t(H+DD^*)}=e^{-tD^*D}D^*.                                \tag{6}
\]
These are exact incidence/gauge identities with the ORIGINAL relative
domains. No arbitrary translation between component grids is used.
Consequently derivative-compatible component changes in (6) have no
overlap-map commutator defect. Source differentiation must retain the
derivatives of $D$ and the complete Hessian, including its curvature
partners; it cannot replace that Hessian by $d_A^*d_A$ alone.

At flat background multiply (4) by its literal propagators:
\[
 \begin{split}
 G_{ia}G_{jb}t_{abc}(GD)_{cv}
 =-&\big[(P_\perp)_{id}G_{jb}A_{dbw}\\
       &+(P_\perp)_{jd}G_{ia}A_{daw}\big]g_{wv},\qquad
 P_\perp=I-DgD^* .                                               \end{split}\tag{7}
\]
This explicitly produces the collapsed local identity contact and its
scalar-gauge correction with the correct signs. The same $A$ occurs
in the ghost-Jacobian vertex $M_1=D^*R_1$.
Equation (5) supplies the quartic and $M_2=D^*R_2$ partners.
The nonlocal part of $P_\perp$ is retained; replacing it by identity
would lose part of the cancellation.

\paragraph{3. The original Haar contact has a definite coefficient.}
In the same coordinates Haar invariance is
$\operatorname{div}\mathcal R_\xi+
 \mathcal R_\xi\cdot\nabla\log j=0$.
Since $\log j=\alpha^2J_2+\cdots$ and
$J_2=-\sum_e|{\bf z}_e|^2/3$ for $SU(2)$,
\[
 \operatorname{div}R_1=0,\qquad
 \boxed{\operatorname{div}(R_2\xi)+
                    \partial_{D\xi}J_2=0.}                      \tag{8}
\]
Explicitly, with either consistent Pauli/quaternion bracket convention,
\[
 (R_2\xi)_e=\tfrac13[{\bf z}_e({\bf z}_e\cdot(D\xi)_e)
                         -|{\bf z}_e|^2(D\xi)_e].
\]
Its three-coordinate divergence is
$2{\bf z}_e\cdot(D\xi)_e/3$, canceled exactly by the Haar derivative.
This is a local coefficient match, not permission to delete Haar from
the full nonlinear bracket. It is also why no additional gauge-parameter
Haar polynomial is added to the left-invariant Jacobian.

\paragraph{4. Exact original-cube checks and scope.}
The finite calculation
expands every original plaquette through fourth order and differentiates
it along the actual fields in (1). On the complete $R=2$ and $R=3$
cubes (216 and 864 plaquettes; 8 and 96 free links), the cubic residual
(2), quartic residual (3), and Haar residual (8) all vanish EXACTLY.
The $R=3$ check includes genuinely nonzero cubic vertices and multiple
interior gauge parameters.

Equations (4)--(8) match the explicit local gauge-sector collapsed-line
contacts to quartic, ghost-Jacobian and Haar partners, with the finite
boundary intact. They explain at tensor level the longitudinal
cancellation of the local noise ward and longitudinal cancellation calculation. They do NOT imply that a curl-contracted
transverse short line is a $D$-contracted leg. The contact isolated in the short line contact before current decay calculation
still requires the complete transverse index sum, its endpoint/spin
vertices and the quartic/ghost/Haar partners, or a
complete boundary subtraction with a controlled marginal kinetic term.
No claim that this entire transverse matching has been completed is
made. The surviving spatial/current estimate and uniform original
remainder remain open.
\end{document}
